\documentclass[11pt,a4paper]{article}
\usepackage[T1]{fontenc}
\usepackage[utf8]{inputenc}
\usepackage{lmodern}
\usepackage[margin=25mm]{geometry}
\usepackage{microtype}
\usepackage[hidelinks]{hyperref}
\hypersetup{
  pdftitle={On completeness conditions on quantale valued sets},
  pdfauthor={Hugo Luiz Mariano},
  pdfsubject={Online Algebra, Logic and Topology Seminar}
}
\setlength{\parindent}{0pt}

\begin{document}
\begin{center}
{\Large\bfseries Online Algebra, Logic and Topology Seminar}\\[4pt]
{\large (CMUC, University of Coimbra)}
\end{center}

\vspace{1em}
\textbf{Speaker:} Hugo Luiz Mariano (Universidade de S\~ao Paulo (IME-USP), Brazil)\\
\textbf{Date:} December 1st, 2026\\
\textbf{Time:} 15:00 (Europe/Lisbon)\\
\textbf{Place:} Google Meet

\vspace{1.5em}
\begin{center}
{\Large\bfseries On completeness conditions on quantale valued sets}\\[6pt]
(joint work with Jos\'e Goudet Alvim and Caio de Andrade Mendes)
\end{center}
\vspace{1em}

\textbf{Abstract.} Categories of quantale valued sets ($Q$-sets) has been proposed over the years in attempts to expand the celebrated notion of $H$-sets, for complete Heyting algebra $H$, to the broader category of parameter algebras consisting of certain quantales---those being (not necessarily commutative) residuated monoids in the category of posets and supremum-preserving maps.

\medskip
In this work, we examine various notions of completeness for $Q$-sets (such as singleton completeness, gluing completeness, and net completeness), where $Q$ is a commutative, semicartesian quantale---a structure that properly extends the notion of a locale and encompasses examples such as the quantale of ideals of a commutative ring with identity, as well as Lawvere's quantale $[0,\infty]$.
\end{document}
